% =====================================================================
%  附录 A · 数学符号表 / Symbol Table
%  按主题分组——配中英文名称 + 章节出处
% =====================================================================

\chapter{数学符号表 / Symbol Table}
\label{app:symbols}

\begin{quote}\itshape
数学家是\textbf{符号的工匠}。\textbf{$\int$ 写不好——读者一辈子记不住}。\\
\textbf{本附录把全书出现过的符号——按主题分门别类——配上中英文名称与首次出现的章节}——\\
\textbf{你忘了哪个符号——回到这里 5 秒钟找到}。
\end{quote}

\section{微积分 / Calculus}

\begin{center}
\renewcommand{\arraystretch}{1.25}
\begin{tabular}{@{}llll@{}}
\toprule
\textbf{符号} & \textbf{中文名称} & \textbf{English Name} & \textbf{章节} \\
\midrule
$\displaystyle \lim_{x \to a} f(x)$ & 极限 & Limit & Ch 1 \\
$\varepsilon, \delta$ & 任意小、足够小 & Epsilon, delta & Ch 1 \\
$f'(x)$ & 一阶导数 & First derivative & Ch 1 \\
$\dfrac{dy}{dx}$ & Leibniz 导数记号 & Leibniz notation & Ch 1 \\
$\dot{x}, \ddot{x}$ & 时间导数（点） & Time derivative (Newton dot) & Ch 1, Ch 4 \\
$f''(x), f^{(n)}(x)$ & 高阶导数 & Higher derivatives & Ch 1 \\
$df$ & 微分 & Differential & Ch 1 \\
$\Delta x$ & 增量 & Increment & Ch 1 \\
$\displaystyle \int f(x)\,dx$ & 不定积分 & Indefinite integral & Ch 1 \\
$\displaystyle \int_a^b f(x)\,dx$ & 定积分 & Definite integral & Ch 1 \\
$\displaystyle \oint$ & 围道积分 & Contour / closed integral & Ch 3, Ch 9 \\
$\sum, \prod$ & 求和、求积 & Sum, product & Ch 1 \\
$e, \pi, \gamma$ & 自然底数、圆周率、Euler 常数 & Math constants & Ch 1 \\
$\infty$ & 无穷 & Infinity & Ch 1 \\
$O(g(n))$ & 大 O 记号 & Big-O notation & Ch 12, Ch 14 \\
$o(g(n))$ & 小 o 记号 & Little-o & Ch 1 \\
\bottomrule
\end{tabular}
\end{center}

\section{多元微积分 / Multivariable Calculus}

\begin{center}
\renewcommand{\arraystretch}{1.25}
\begin{tabular}{@{}llll@{}}
\toprule
\textbf{符号} & \textbf{中文名称} & \textbf{English Name} & \textbf{章节} \\
\midrule
$\dfrac{\partial f}{\partial x}$ & 偏导数 & Partial derivative & Ch 2 \\
$\nabla f$ & 梯度 & Gradient & Ch 2, Ch 13 \\
$\nabla^2 f$ 或 $\Delta f$ & Laplace 算子 & Laplacian & Ch 3, Ch 5 \\
$\nabla \cdot \mathbf{F}$ & 散度 & Divergence & Ch 3 \\
$\nabla \times \mathbf{F}$ & 旋度 & Curl & Ch 3 \\
$\mathbf{H} = \nabla^2 f$ & Hessian 矩阵 & Hessian matrix & Ch 2, Ch 13 \\
$\mathbf{J}$ & Jacobian 矩阵 & Jacobian matrix & Ch 2 \\
$J = \det \mathbf{J}$ & Jacobian 行列式 & Jacobian determinant & Ch 2 \\
$d\mathbf{r}, d\mathbf{S}, dV$ & 弧元、面元、体元 & Line/area/volume element & Ch 3 \\
$\iint, \iiint$ & 二重、三重积分 & Double, triple integral & Ch 2 \\
\bottomrule
\end{tabular}
\end{center}

\section{矢量分析 / Vector Analysis}

\begin{center}
\renewcommand{\arraystretch}{1.25}
\begin{tabular}{@{}llll@{}}
\toprule
\textbf{符号} & \textbf{中文名称} & \textbf{English Name} & \textbf{章节} \\
\midrule
$\mathbf{v}, \vec{v}$ & 矢量 & Vector & Ch 3, Ch 6 \\
$\hat{\mathbf{n}}$ & 单位法矢 & Unit normal & Ch 3 \\
$\mathbf{a} \cdot \mathbf{b}$ & 点积 / 内积 & Dot product / inner product & Ch 3, Ch 6 \\
$\mathbf{a} \times \mathbf{b}$ & 叉积 & Cross product & Ch 3 \\
$\|\mathbf{v}\|$ 或 $|\mathbf{v}|$ & 矢量模长 & Norm / magnitude & Ch 3, Ch 6 \\
$\mathbf{\hat{e}}_x, \mathbf{\hat{e}}_y, \mathbf{\hat{e}}_z$ & 直角坐标基矢 & Cartesian basis & Ch 3 \\
$\mathbf{\hat{e}}_r, \mathbf{\hat{e}}_\theta, \mathbf{\hat{e}}_\phi$ & 球坐标基矢 & Spherical basis & Ch 3 \\
$\displaystyle \oint_{\partial \Omega}$ & 边界曲线/曲面积分 & Boundary integral & Ch 3 \\
$\mathbf{E}, \mathbf{B}$ & 电场、磁场 & Electric, magnetic field & Ch 3 \\
\bottomrule
\end{tabular}
\end{center}

\section{线性代数 / Linear Algebra}

\begin{center}
\renewcommand{\arraystretch}{1.25}
\begin{tabular}{@{}llll@{}}
\toprule
\textbf{符号} & \textbf{中文名称} & \textbf{English Name} & \textbf{章节} \\
\midrule
$\mathbf{A}, A_{ij}$ & 矩阵及其元素 & Matrix and entries & Ch 6, Ch 7 \\
$\mathbf{A}^T$ & 转置 & Transpose & Ch 6 \\
$\mathbf{A}^{-1}$ & 逆矩阵 & Inverse & Ch 6 \\
$\mathbf{A}^\dagger$ & 共轭转置 & Conjugate transpose (Hermitian) & Ch 7, Ch 9 \\
$\mathbf{A}^+$ & 伪逆（Moore-Penrose）& Pseudoinverse & Ch 8 \\
$\det \mathbf{A}, |\mathbf{A}|$ & 行列式 & Determinant & Ch 6 \\
$\operatorname{tr}(\mathbf{A})$ & 迹 & Trace & Ch 7 \\
$\operatorname{rank}(\mathbf{A})$ & 秩 & Rank & Ch 6, Ch 8 \\
$\ker \mathbf{A}, \operatorname{null}$ & 零空间 / 核 & Null space / kernel & Ch 8 \\
$\operatorname{im} \mathbf{A}, \operatorname{col}$ & 列空间 / 像 & Column space / image & Ch 8 \\
$\mathbf{I}$ & 单位矩阵 & Identity & Ch 6 \\
$\mathbf{0}$ & 零矩阵 / 零矢量 & Zero matrix / vector & Ch 6 \\
$\lambda$ & 特征值 & Eigenvalue & Ch 7 \\
$\mathbf{v}_\lambda$ & 特征矢量 & Eigenvector & Ch 7 \\
$\mathbf{A} = \mathbf{U} \boldsymbol{\Sigma} \mathbf{V}^T$ & 奇异值分解 SVD & Singular Value Decomposition & Ch 7, Ch 8 \\
$\sigma_i$ & 奇异值 & Singular value & Ch 7 \\
$\|\mathbf{A}\|_2, \|\mathbf{A}\|_F$ & 谱范数、Frobenius 范数 & Spectral / Frobenius norm & Ch 7 \\
$\langle \mathbf{u}, \mathbf{v} \rangle$ & 内积 & Inner product & Ch 6 \\
$\mathbf{u} \otimes \mathbf{v}$ & 外积 / 张量积 & Outer / tensor product & Ch 7 \\
$\oplus, \perp$ & 直和、正交 & Direct sum, orthogonal & Ch 8 \\
\bottomrule
\end{tabular}
\end{center}

\section{微分方程 ODE / PDE}

\begin{center}
\renewcommand{\arraystretch}{1.25}
\begin{tabular}{@{}llll@{}}
\toprule
\textbf{符号} & \textbf{中文名称} & \textbf{English Name} & \textbf{章节} \\
\midrule
$y' + p(x)y = q(x)$ & 一阶线性 ODE & First-order linear ODE & Ch 4 \\
$ay'' + by' + cy = f(x)$ & 二阶线性 ODE & Second-order linear ODE & Ch 4 \\
$\dot{\mathbf{x}} = \mathbf{A}\mathbf{x}$ & 状态空间 ODE & State-space ODE & Ch 4 \\
$\omega_0, \zeta$ & 自然频率、阻尼比 & Natural freq, damping ratio & Ch 4 \\
$\dfrac{\partial u}{\partial t} = \alpha \nabla^2 u$ & 热传导方程 & Heat equation & Ch 5 \\
$\dfrac{\partial^2 u}{\partial t^2} = c^2 \nabla^2 u$ & 波动方程 & Wave equation & Ch 5 \\
$\nabla^2 u = 0$ & Laplace 方程 & Laplace equation & Ch 5 \\
$\nabla^2 u = f$ & Poisson 方程 & Poisson equation & Ch 5 \\
$u_t + u u_x = \nu u_{xx}$ & Burgers 方程 & Burgers equation & Ch 5 \\
$u|_{\partial \Omega}, \partial u / \partial n|_{\partial \Omega}$ & Dirichlet / Neumann 边界 & Dirichlet / Neumann BCs & Ch 5 \\
$G(\mathbf{r}, \mathbf{r}')$ & Green 函数 & Green's function & Ch 5 \\
\bottomrule
\end{tabular}
\end{center}

\section{复数与复分析 / Complex Numbers and Analysis}

\begin{center}
\renewcommand{\arraystretch}{1.25}
\begin{tabular}{@{}llll@{}}
\toprule
\textbf{符号} & \textbf{中文名称} & \textbf{English Name} & \textbf{章节} \\
\midrule
$i, j$ & 虚单位 & Imaginary unit & Ch 9 \\
$z = a + bi$ & 复数 & Complex number & Ch 9 \\
$\bar{z}, z^*$ & 共轭复数 & Complex conjugate & Ch 9 \\
$|z|, \arg z$ & 模、辐角 & Modulus, argument & Ch 9 \\
$\operatorname{Re}(z), \operatorname{Im}(z)$ & 实部、虚部 & Real, imaginary parts & Ch 9 \\
$e^{i\theta} = \cos\theta + i\sin\theta$ & Euler 公式 & Euler's formula & Ch 9 \\
$f: \mathbb{C} \to \mathbb{C}$ & 复函数 & Complex function & Ch 9 \\
$\displaystyle \oint_C f(z)\,dz$ & 复围道积分 & Contour integral & Ch 9 \\
$\operatorname{Res}_{z = z_0} f$ & 留数 & Residue & Ch 9 \\
$\mathcal{L}\{f\}(s)$ & Laplace 变换 & Laplace transform & Ch 4, Ch 9 \\
$\mathcal{F}\{f\}(\omega), \hat{f}$ & Fourier 变换 & Fourier transform & Ch 5, Ch 9 \\
$\mathcal{Z}\{x_n\}(z)$ & Z 变换 & Z-transform & Ch 9 \\
$H(s), H(j\omega), H(z)$ & 传递函数 & Transfer function & Ch 4, Ch 9 \\
\bottomrule
\end{tabular}
\end{center}

\section{概率与统计 / Probability and Statistics}

\begin{center}
\renewcommand{\arraystretch}{1.25}
\begin{tabular}{@{}llll@{}}
\toprule
\textbf{符号} & \textbf{中文名称} & \textbf{English Name} & \textbf{章节} \\
\midrule
$\Omega, \mathcal{F}, P$ & 样本空间、$\sigma$-代数、概率测度 & Sample space, sigma-algebra, measure & Ch 10 \\
$P(A)$ & 事件 $A$ 的概率 & Probability of event $A$ & Ch 10 \\
$P(A \mid B)$ & 条件概率 & Conditional probability & Ch 10 \\
$X, Y$ & 随机变量 & Random variable & Ch 10 \\
$\mathbb{E}[X], \mu$ & 期望、均值 & Expectation, mean & Ch 10 \\
$\operatorname{Var}(X), \sigma^2$ & 方差 & Variance & Ch 10 \\
$\sigma$ & 标准差 & Standard deviation & Ch 10 \\
$\operatorname{Cov}(X,Y), \rho$ & 协方差、相关系数 & Covariance, correlation & Ch 10 \\
$f_X(x), F_X(x)$ & PDF、CDF & Density / distribution function & Ch 10 \\
$X \sim \mathcal{N}(\mu, \sigma^2)$ & 正态分布 & Normal distribution & Ch 10 \\
$X \sim \mathrm{Bin}(n,p)$ & 二项分布 & Binomial distribution & Ch 10 \\
$X \sim \mathrm{Poisson}(\lambda)$ & Poisson 分布 & Poisson distribution & Ch 10 \\
$X \sim \mathrm{Exp}(\lambda)$ & 指数分布 & Exponential distribution & Ch 10 \\
$\bar{x}, s^2$ & 样本均值、样本方差 & Sample mean / variance & Ch 11 \\
$H_0, H_1$ & 原假设、备择假设 & Null / alternative hypothesis & Ch 11 \\
$p$ & $p$ 值 & $p$-value & Ch 11 \\
$\alpha, \beta$ & 显著性水平、第二类错误率 & Significance, type-II error rate & Ch 11 \\
$\hat{\theta}$ & 参数估计 & Parameter estimate & Ch 11 \\
$L(\theta), \ell(\theta)$ & 似然、对数似然 & Likelihood, log-likelihood & Ch 11 \\
$\operatorname{KL}(p \| q)$ & KL 散度 & KL divergence & Ch 11 \\
$\chi^2, t, F$ & 卡方、$t$、$F$ 分布 & Chi-squared, $t$, $F$ distributions & Ch 11 \\
\bottomrule
\end{tabular}
\end{center}

\section{数值与优化 / Numerical and Optimization}

\begin{center}
\renewcommand{\arraystretch}{1.25}
\begin{tabular}{@{}llll@{}}
\toprule
\textbf{符号} & \textbf{中文名称} & \textbf{English Name} & \textbf{章节} \\
\midrule
$\kappa(\mathbf{A})$ & 条件数 & Condition number & Ch 12 \\
$\varepsilon_{\text{mach}}$ & 机器精度 & Machine epsilon & Ch 12 \\
$h$ & 步长 & Step size & Ch 12 \\
$\arg\min, \arg\max$ & 极小/极大化变量 & Argmin / argmax & Ch 13 \\
$\min, \max$ & 极小/极大值 & Min / max & Ch 13 \\
$\nabla f, \mathbf{H}$ & 梯度、Hessian & Gradient, Hessian & Ch 13 \\
$\mathcal{L}(\mathbf{x}, \boldsymbol{\lambda})$ & Lagrangian 函数 & Lagrangian & Ch 13 \\
$\boldsymbol{\lambda}, \boldsymbol{\mu}$ & 对偶变量 & Dual variables (KKT multipliers) & Ch 13 \\
$\eta, \alpha$ & 学习率 & Learning rate & Ch 13 \\
$\beta_1, \beta_2$ & Adam 动量参数 & Adam momentum parameters & Ch 13 \\
\bottomrule
\end{tabular}
\end{center}

\section{集合、逻辑与离散 / Sets, Logic, Discrete (Ch 14)}

\begin{center}
\renewcommand{\arraystretch}{1.25}
\begin{tabular}{@{}llll@{}}
\toprule
\textbf{符号} & \textbf{中文名称} & \textbf{English Name} & \textbf{章节} \\
\midrule
$\mathbb{N}, \mathbb{Z}, \mathbb{Q}, \mathbb{R}, \mathbb{C}$ & 自然/整/有理/实/复数集 & Number sets & 通用 \\
$\in, \notin, \subset, \subseteq$ & 属于、包含 & Belongs to, subset & 通用 \\
$\cup, \cap, \setminus$ & 并、交、差 & Union, intersection, difference & Ch 10, Ch 14 \\
$\emptyset$ & 空集 & Empty set & 通用 \\
$|A|$ & 集合基数 & Cardinality & Ch 14 \\
$\binom{n}{k}, C(n,k)$ & 二项系数 & Binomial coefficient & Ch 10, Ch 14 \\
$P(n,k)$ & 排列数 & Number of permutations & Ch 14 \\
$n!$ & 阶乘 & Factorial & Ch 14 \\
$\land, \lor, \neg$ & 与、或、非 & AND, OR, NOT & Ch 14 \\
$\to, \leftrightarrow$ & 蕴含、等价 & Implies, iff & Ch 14 \\
$\forall, \exists$ & 任意、存在 & For all, exists & 通用 \\
$\equiv$ & 恒等 / 同余 & Equivalent / congruent & Ch 14 \\
$G = (V, E)$ & 图（点集 + 边集） & Graph & Ch 14 \\
$\deg(v)$ & 顶点度数 & Vertex degree & Ch 14 \\
$O, \Omega, \Theta$ & 复杂度三符号 & Asymptotic notation & Ch 14 \\
$P, NP, NP\text{-complete}$ & 复杂度类 & Complexity classes & Ch 14 \\
\bottomrule
\end{tabular}
\end{center}

\section{常用希腊字母 / Greek Letters}

\begin{center}
\renewcommand{\arraystretch}{1.2}
\begin{tabular}{@{}llll@{}}
\toprule
\textbf{符号} & \textbf{名称} & \textbf{常见用法} & \textbf{章节} \\
\midrule
$\alpha$ & alpha & 学习率、显著性水平、热扩散率 & Ch 5, Ch 11, Ch 13 \\
$\beta$ & beta & 阻尼、Adam 动量、回归系数 & Ch 4, Ch 11, Ch 13 \\
$\gamma$ & gamma & Euler 常数、折扣因子 & Ch 1 \\
$\delta$ & delta & 任意小、Dirac $\delta$ 函数 & Ch 1, Ch 5 \\
$\varepsilon$ & epsilon & 误差容限、机器精度 & Ch 1, Ch 12 \\
$\zeta$ & zeta & 阻尼比、Riemann zeta & Ch 4 \\
$\eta$ & eta & 学习率、效率 & Ch 13 \\
$\theta, \vartheta$ & theta & 角度、参数 & Ch 3, Ch 11 \\
$\kappa$ & kappa & 曲率、条件数 & Ch 12 \\
$\lambda$ & lambda & 特征值、Lagrange 乘子、波长 & Ch 7, Ch 13 \\
$\mu$ & mu & 均值、磁导率、KKT 乘子 & Ch 10, Ch 13 \\
$\nu$ & nu & 粘度、自由度 & Ch 5 \\
$\xi$ & xi & 中间点、辅助变量 & Ch 1 \\
$\pi$ & pi & 圆周率 & 通用 \\
$\rho$ & rho & 密度、相关系数 & Ch 5, Ch 10 \\
$\sigma$ & sigma & 标准差、奇异值、$\sigma$-代数 & Ch 7, Ch 10 \\
$\tau$ & tau & 时间常数、剪应力 & Ch 4 \\
$\phi, \varphi$ & phi & 相位、标量势 & Ch 3, Ch 9 \\
$\chi$ & chi & 卡方分布、电极化率 & Ch 11 \\
$\psi$ & psi & 波函数、流函数 & Ch 5 \\
$\omega$ & omega & 角频率 & Ch 4, Ch 9 \\
$\Delta$ & Delta & 增量、Laplace 算子 & Ch 1, Ch 5 \\
$\Sigma$ & Sigma & 求和符号、协方差矩阵 & Ch 1, Ch 10 \\
$\Pi$ & Pi & 求积符号 & Ch 10 \\
$\Phi$ & Phi & 标量势、累积正态分布 & Ch 3, Ch 10 \\
$\Omega$ & Omega & 区域、样本空间、阻抗 & Ch 5, Ch 10 \\
\bottomrule
\end{tabular}
\end{center}

\section{使用建议 / How to Use This Table}

\begin{enumerate}
\item \textbf{读外文论文}——\textbf{遇到不认识的符号——\textbf{先查\textbf{对应表}}}——\textbf{再回到正文章节查\textbf{定义与直觉}}。
\item \textbf{写论文 / 写代码注释}——\textbf{用本表保持\textbf{符号一致性}}——\textbf{别同一篇里 $\sigma$ 既表标准差又表奇异值}。
\item \textbf{讲课 / 写文档}——\textbf{先列\textbf{符号表}}（"Notation"）——\textbf{是学术写作的\textbf{专业素养}}。
\item \textbf{看物理 vs 数学教材}——\textbf{注意\textbf{同一符号在不同领域含义不同}}：例如 $j$ 在数学里是虚单位、在电气工程里也是虚单位（避开电流 $i$）、在物理里是角动量量子数。
\end{enumerate}

\clearpage
